\section{Venice, England, and the Reach of Indirect Power}

\subsection{The distinction becomes a political instrument}

Bellarmine's denial of direct universal papal temporal dominion had once drawn
the anger of Sixtus~V.  It did not withdraw public affairs from papal judgment.
His indirect-power theory allowed intervention where temporal decisions bore
on the spiritual good of the Church.  The limit was real within contemporary
Catholic debate, but so was the reach that remained.

Motta's inspected critical reconstruction places Bellarmine among the Roman
writers against Venetian theologians and Paolo Sarpi during the Interdict
controversy of 1606.  The exact Venetian writings were not collated here.
Questions of jurisdiction, obedience, and the relation of civil law to
ecclesiastical authority were being contested by a republic and the papacy,
not merely by authors in a classroom.  A full account requires Bellarmine's
texts, Venetian state records, and Sarpi's texts in their own setting; this
biography does not let one side's controversial categories define the other
without remainder.

The English Oath of Allegiance generated another sequence from 1606 into 1609.
Bellarmine corresponded with George Blackwell, answered James~I, and published
under the name Matteo Torti.  He did not travel to England for a personal
diplomatic negotiation.  His engagement was principally textual and curial,
but its effects were not merely literary.  English Catholics had to negotiate
loyalty to a monarch, communion with Rome, legal penalties, and the risks
created by arguments made at a safe distance.

In 1610 Bellarmine published \emph{De potestate Summi Pontificis in rebus
temporalibus} against William Barclay.  The work restated the disputed boundary
under new political pressure.  Denial of direct dominion did not mean that a
ruler's temporal acts were immune from papal consequence when Bellarmine judged
spiritual goods at stake.  For that reason he should neither be presented as a
theorist of unlimited direct theocracy nor recruited as a straightforward
ancestor of modern liberal separation.

\subsection{Argument, conscience, and subjects at risk}

Bellarmine's political theology was ecclesiological before it was
constitutional in a modern sense.  It asked how the Church's supernatural end
related to the powers of rulers and the obligations of baptized persons.  Its
distinctions could restrain a claim of immediate temporal ownership while
authorizing serious intervention through spiritual jurisdiction.  Whether a
specific intervention was prudent or just could not be decided merely by
labeling the power ``indirect.''

The Venetian and English controversies also expose the asymmetry of learned
debate.  Authors could refine jurisdictional claims while subjects bore the
legal and social risks of disputed allegiance.  A complete moral history would
need their evidence, including Catholic voices not represented by Bellarmine's
strategy.  The current source base controls his publications and the broad
sequence, not the whole reciprocal dossier or every practical effect.

These boundaries do not render his distinctions meaningless.  They identify a
serious attempt to limit one account of papal temporal power and to preserve a
spiritual logic for intervention.  Historical precision permits both
recognitions: the theory was contested, bounded, and intellectually important;
it also remained confessional, coercion-compatible, and distant from present
assumptions about equal citizenship and religious liberty.

\evidenceaxis{Bellarmine denied direct universal papal temporal dominion but
defended indirect intervention for spiritual ends.}{\emph{De Romano
Pontifice} V, especially chapters 6--7, and \emph{De potestate} (1610) control
the theory; Motta controls the broad Venetian and oath sequence because the
exact 1606 Venetian writings were not collated.}{Indirect power is not modern
church--state separation; opponents, vulnerable subjects, and the practical
results require reciprocal political and legal records.}
